% =============================================================================
% APPENDIX B: CORE-HOW: Complementarity and Interaction Effects
% =============================================================================
% Module: BCM2_02_CMP (Complementarity)
% Category: CORE
% Version: 1.0 (January 2026)
% Status: Final
% Language: English
%
% PURPOSE: This appendix answers the fundamental question "How do utility
% components interact?" It formalizes the Complementarity Parameter γ,
% establishing that utility dimensions and levels can be complements (γ > 0),
% substitutes (γ < 0), or independent (γ = 0).
%
% 10C POSITION: This is CORE 3 of 10 - the "HOW" question.
%
% =============================================================================

% -----------------------------------------------------------------------------
% HEADER BLOCK
% -----------------------------------------------------------------------------

\begin{tcolorbox}[colback=gray!5!white, colframe=gray!75!black,
    title=Appendix Header: B CORE-HOW]

\begin{tabular}{@{}ll@{}}
\textbf{Code:} & B \\
\textbf{Category:} & CORE (Fundamental Theory) \\
\textbf{Full Name:} & CORE-HOW: Complementarity and Interaction Effects \\
\textbf{Module:} & BCM2\_02\_CMP (Complementarity) \\
\textbf{Version:} & 1.0 (January 2026) \\
\textbf{Status:} & Final \\
\textbf{Language:} & English \\
\end{tabular}

\vspace{0.5em}
\textbf{Purpose:} Formalizes \textbf{Complementarity Theory}, answering the
third fundamental question of the EBF: \textit{``How do utility components interact?''}
Establishes that interactions are characterized by the complementarity parameter
$\gamma_{ij} \in [-1, +1]$, where positive values indicate synergy (complements),
negative values indicate crowding-out (substitutes), and zero indicates independence.

\textbf{10C Position:} This is CORE 3 (HOW) --- builds on WHO (AAA) and WHAT (C),
provides interaction structure for all EBF models.

\end{tcolorbox}

% =============================================================================
% CROSS-REFERENCE STRUCTURE
% =============================================================================
%
% CHAPTER LINKAGE:
% - Primary Chapter: Chapter 7: Complementarity and Non-Concavity
% - Secondary Chapters: Ch. 8 (Mathematical Foundations), Ch. 17-20 (Interventions)
%
% DEPENDENCIES:
% - Appendix UNMAPPED_AAA (CORE-WHO): Level structure for cross-level complementarity
% - Appendix UNMAPPED_C (CORE-WHAT): Dimension structure for cross-dimension complementarity
%
% DEPENDENTS:
% - Appendix UNMAPPED_V (CORE-WHEN): Context-dependent γ values
% - Appendix UNMAPPED_IE (CORE-EIT): Intervention complementarity/crowding
% - Appendix UNMAPPED_CAL (METHOD-LLMMC): γ calibration methods
%
% =============================================================================

\begin{tcolorbox}[colback=purple!5!white, colframe=purple!75!black,
    title=Chapter Linkage]

\textbf{Primary Chapter:} Chapter 7: Complementarity and Non-Concavity

\vspace{0.3em}
\textbf{The Fundamental Question:}
\begin{quote}
\textit{``How do utility components interact? When do they enhance each other (complements) and when do they undermine each other (substitutes)?''}
\end{quote}

\textbf{Section-Level Mapping:}
\begin{itemize}[nosep]
\item Section 7.1: Why Interaction Matters $\leftrightarrow$ Section~\ref{sec:how-fundamental}
\item Section 7.2: The γ Parameter $\leftrightarrow$ Section~\ref{sec:how-gamma}
\item Section 7.3: Cross-Dimension Complementarity $\leftrightarrow$ Section~\ref{sec:how-dimension}
\item Section 7.4: Cross-Level Complementarity $\leftrightarrow$ Section~\ref{sec:how-level}
\end{itemize}

\textbf{Secondary Chapters:}
\begin{itemize}[nosep]
\item Chapter 8: Mathematical Foundations --- Non-concavity and emergence
\item Chapter 17: Intervention Foundations --- Crowding-out in intervention portfolios
\item Chapter 20: Intervention Portfolios --- γ-matrix for portfolio optimization
\end{itemize}

\textbf{Relationship to Other COREs:}
\begin{itemize}[nosep]
\item Appendix UNMAPPED_AAA (CORE-WHO): Provides levels whose interactions HOW describes
\item Appendix UNMAPPED_C (CORE-WHAT): Provides dimensions whose interactions HOW describes
\item Appendix UNMAPPED_V (CORE-WHEN): Provides context that modulates γ values
\end{itemize}

\end{tcolorbox}

% -----------------------------------------------------------------------------
% APPENDIX SCOPE BOX
% -----------------------------------------------------------------------------

\begin{tcolorbox}[
    colback=orange!5!white,
    colframe=orange!75!black,
    title={\textbf{Appendix Scope: Ziel / In-Scope / Out-of-Scope / Constraints}},
    fonttitle=\bfseries
]

\textbf{Ziel (Objective):}
\begin{itemize}[nosep]
    \item Define how utility components interact through the complementarity parameter γ
\end{itemize}

\textbf{In-Scope:}
\begin{itemize}[nosep]
    \item Definition and properties of γ (axioms UNMAPPED_CMP-1 through UNMAPPED_CMP-10)
    \item Cross-dimension complementarity (γ between FEPSDE dimensions)
    \item Cross-level complementarity (γ between welfare levels)
    \item Implications for non-concavity and emergence
\end{itemize}

\textbf{Out-of-Scope:}
\begin{itemize}[nosep]
    \item What the dimensions are $\rightarrow$ Appendix UNMAPPED_C (CORE-WHAT)
    \item What the levels are $\rightarrow$ Appendix UNMAPPED_AAA (CORE-WHO)
    \item When γ varies with context $\rightarrow$ Appendix UNMAPPED_V (CORE-WHEN)
    \item How to estimate γ empirically $\rightarrow$ Appendix UNMAPPED_BBB (CORE-WHERE)
\end{itemize}

\textbf{Constraints:}
\begin{itemize}[nosep]
    \item γ is bounded: $\gamma_{ij} \in [-1, +1]$
    \item γ is symmetric by default: $\gamma_{ij} = \gamma_{ji}$
    \item γ = 0 implies separability (standard economics assumption)
\end{itemize}

\textbf{Lieferobjekte (Deliverables):}
\begin{itemize}[nosep]
    \item Complete γ taxonomy (Table~\ref{tab:how-gamma-types})
    \item Axiom System UNMAPPED_CMP-1 through UNMAPPED_CMP-10
    \item Cross-Dimension γ Matrix template
    \item Cross-Level γ Matrix template
\end{itemize}

\end{tcolorbox}

% =============================================================================
% CHAPTER
% =============================================================================

\chapter{CORE-HOW: Complementarity and Interaction Effects}
\label{app:core-how}

\begin{abstract}
This appendix answers the third fundamental question of the Evidence-Based Framework:
\textit{``How do utility components interact?''} Standard economics assumes additive
separability --- components contribute independently. We reject this assumption and
formalize the complementarity parameter $\gamma_{ij} \in [-1, +1]$, capturing synergies
($\gamma > 0$) and crowding-out ($\gamma < 0$) between utility dimensions and welfare
levels. The axiom system UNMAPPED_CMP-1 through UNMAPPED_CMP-10 provides the formal foundations.
This interaction structure is essential for understanding non-concave utility,
intervention portfolio design, and emergence phenomena.
\end{abstract}

% -----------------------------------------------------------------------------
% QUICK REFERENCE BOX
% -----------------------------------------------------------------------------

\begin{tcolorbox}[colback=blue!5!white, colframe=blue!75!black,
    title=Quick Reference: CORE-HOW]

\textbf{Core Question:} How do utility components interact?

\textbf{Key Formula:}
\begin{equation}
U = \sum_i u_i + \sum_{i < j} \gamma_{ij} \cdot u_i \cdot u_j
\label{eq:how-utility}
\end{equation}

\textbf{Key Concepts:}
\begin{itemize}[nosep]
\item \textbf{Complementarity $(\gamma > 0)$}: Components enhance each other
\item \textbf{Substitutability $(\gamma < 0)$}: Components crowd each other out
\item \textbf{Separability $(\gamma = 0)$}: Components are independent
\item \textbf{γ-Matrix}: Complete specification of all pairwise interactions
\end{itemize}

\textbf{Cross-References:}
AAA (CORE-WHO), C (CORE-WHAT), V (CORE-WHEN), IE (CORE-EIT), CAL (METHOD-LLMMC), CP (DOMAIN-CPS: Vernetztheit γ)

\end{tcolorbox}


%%%%%%%%%%%%%%%%%%%%%%%%%%%%%%%%%%%%%%%%%%%%%%%%%%%%%%%%%%%%%%%%%%%%%%%%%%%%%%%
\section{The Fundamental Question}
\label{sec:how-fundamental}
%%%%%%%%%%%%%%%%%%%%%%%%%%%%%%%%%%%%%%%%%%%%%%%%%%%%%%%%%%%%%%%%%%%%%%%%%%%%%%%

\subsection{Why Interaction Matters}

Standard economic models assume utility is additively separable:
\begin{equation}
U = u_F + u_E + u_P + u_S + u_D + u_X
\end{equation}

This implies each dimension contributes independently. But reality is different:

\begin{itemize}
\item Financial security ($u_F$) and social belonging ($u_S$) may be complements:
      having both is more than the sum of having each alone
\item Extrinsic rewards ($u_F$) and intrinsic motivation ($u_X$) may be substitutes:
      adding financial incentives can undermine intrinsic motivation
\item Autonomy and relatedness may interact non-linearly depending on personality
\end{itemize}

\subsection{The Core Problem}

Ignoring interactions leads to:
\begin{enumerate}
\item \textbf{Suboptimal Intervention Design}: Treating components as independent
      misses synergies and ignores crowding-out
\item \textbf{Incorrect Predictions}: Additive models predict linear effects
      when reality shows non-linear emergence
\item \textbf{Policy Failures}: Interventions that look good in isolation
      may fail in combination (or succeed unexpectedly)
\end{enumerate}

\begin{tcolorbox}[colback=red!5!white, colframe=red!75!black,
    title=The Fundamental Question]
\textbf{How do utility components interact? When are they complements (enhancing each other) and when are they substitutes (crowding each other out)?}
\end{tcolorbox}


%%%%%%%%%%%%%%%%%%%%%%%%%%%%%%%%%%%%%%%%%%%%%%%%%%%%%%%%%%%%%%%%%%%%%%%%%%%%%%%
\section{Core Theory: The Complementarity Parameter}
\label{sec:how-gamma}
%%%%%%%%%%%%%%%%%%%%%%%%%%%%%%%%%%%%%%%%%%%%%%%%%%%%%%%%%%%%%%%%%%%%%%%%%%%%%%%

\subsection{Definition}

\begin{definition}[Complementarity Parameter]
The \textbf{complementarity parameter} $\gamma_{ij}$ measures the interaction
between utility components $i$ and $j$:
\begin{equation}
\gamma_{ij} = \frac{\partial^2 U}{\partial u_i \partial u_j}
\end{equation}
normalized to the interval $[-1, +1]$.
\end{definition}

\subsection{Interpretation}

\begin{table}[htbp]
\centering
\caption{Interpretation of γ Values}
\label{tab:how-gamma-types}
\renewcommand{\arraystretch}{1.3}
\begin{tabular}{@{}ccl@{}}
\toprule
\textbf{γ Value} & \textbf{Type} & \textbf{Interpretation} \\
\midrule
$\gamma = +1$ & Perfect Complements & Components must be present together \\
$0 < \gamma < 1$ & Complements & Components enhance each other \\
$\gamma = 0$ & Independent & No interaction (additive separability) \\
$-1 < \gamma < 0$ & Substitutes & Components crowd each other out \\
$\gamma = -1$ & Perfect Substitutes & Components are redundant \\
\bottomrule
\end{tabular}
\end{table}

\subsection{The γ-Matrix}

For $n$ components, the complete interaction structure is captured by:

\begin{equation}
\mathbf{\Gamma} = \begin{pmatrix}
1 & \gamma_{12} & \gamma_{13} & \cdots & \gamma_{1n} \\
\gamma_{21} & 1 & \gamma_{23} & \cdots & \gamma_{2n} \\
\gamma_{31} & \gamma_{32} & 1 & \cdots & \gamma_{3n} \\
\vdots & \vdots & \vdots & \ddots & \vdots \\
\gamma_{n1} & \gamma_{n2} & \gamma_{n3} & \cdots & 1
\end{pmatrix}
\label{eq:how-gamma-matrix}
\end{equation}

with diagonal elements $\gamma_{ii} = 1$ and symmetry $\gamma_{ij} = \gamma_{ji}$ by default.


%%%%%%%%%%%%%%%%%%%%%%%%%%%%%%%%%%%%%%%%%%%%%%%%%%%%%%%%%%%%%%%%%%%%%%%%%%%%%%%
\section{Axioms and Formal Definitions}
\label{sec:how-axioms}
%%%%%%%%%%%%%%%%%%%%%%%%%%%%%%%%%%%%%%%%%%%%%%%%%%%%%%%%%%%%%%%%%%%%%%%%%%%%%%%

% UNMAPPED_CMP-1
\begin{tcolorbox}[colback=gray!5!white, colframe=gray!75!black,
    title=Axiom UNMAPPED_CMP-1: Existence of Interaction]
\textbf{Statement:}
\begin{equation}
\exists \; i, j: \gamma_{ij} \neq 0
\end{equation}

\textbf{Interpretation:} At least some utility components interact non-trivially.
The world is not additively separable.

\textbf{Implication:} Standard economics (which assumes $\gamma_{ij} = 0 \; \forall i,j$)
is a special case, not the general case.

\textbf{Testable Prediction:} Experimental manipulation of multiple components
simultaneously will show interaction effects beyond additive predictions.
\end{tcolorbox}

% UNMAPPED_CMP-2
\begin{tcolorbox}[colback=gray!5!white, colframe=gray!75!black,
    title=Axiom UNMAPPED_CMP-2: Boundedness]
\textbf{Statement:}
\begin{equation}
\forall \; i, j: \gamma_{ij} \in [-1, +1]
\end{equation}

\textbf{Interpretation:} Interactions are bounded. No infinite synergy or infinite crowding.

\textbf{Implication:} The γ-matrix has finite spectral norm.
\end{tcolorbox}

% UNMAPPED_CMP-3
\begin{tcolorbox}[colback=gray!5!white, colframe=gray!75!black,
    title=Axiom UNMAPPED_CMP-3: Default Symmetry]
\textbf{Statement:}
\begin{equation}
\gamma_{ij} = \gamma_{ji} \quad \text{by default}
\end{equation}

\textbf{Interpretation:} The effect of $i$ on the marginal utility of $j$ equals
the effect of $j$ on the marginal utility of $i$.

\textbf{Exception:} Asymmetry can be introduced when theoretically justified
(e.g., temporal precedence, power asymmetries).

\textbf{Cross-Reference:} Asymmetric cases handled in CORE-WHEN (V) via context.
\end{tcolorbox}

% UNMAPPED_CMP-4
\begin{tcolorbox}[colback=gray!5!white, colframe=gray!75!black,
    title=Axiom UNMAPPED_CMP-4: Context Dependence]
\textbf{Statement:}
\begin{equation}
\gamma_{ij} = \gamma_{ij}(\Psi) \quad \text{where } \Psi \text{ is the context vector}
\end{equation}

\textbf{Interpretation:} Complementarity is not fixed; it varies with context.
Components that are complements in one context may be substitutes in another.

\textbf{Example:} Financial incentives and intrinsic motivation:
\begin{itemize}
\item In commercial contexts: often substitutes ($\gamma < 0$)
\item In prosocial contexts: sometimes complements ($\gamma > 0$)
\end{itemize}

\textbf{Cross-Reference:} Full specification of $\gamma(\Psi)$ in CORE-WHEN (V).
\end{tcolorbox}

% UNMAPPED_CMP-5
\begin{tcolorbox}[colback=gray!5!white, colframe=gray!75!black,
    title=Axiom UNMAPPED_CMP-5: Crowding-Out Theorem]
\textbf{Statement:}
\begin{equation}
\gamma_{FX} < 0 \quad \text{(extrinsic crowds out intrinsic, ceteris paribus)}
\end{equation}

\textbf{Interpretation:} Financial incentives ($u_F$) and intrinsic motivation ($u_X$)
are typically substitutes, not complements. This is the classic crowding-out effect.

\textbf{Empirical Basis:} \citet{frey1997relationship}, \citet{deci1999meta}

\textbf{Boundary Conditions:} See UNMAPPED_CMP-4 --- the sign can flip in specific contexts.
\end{tcolorbox}

% UNMAPPED_CMP-6
\begin{tcolorbox}[colback=gray!5!white, colframe=gray!75!black,
    title=Axiom UNMAPPED_CMP-6: Social-Financial Tension]
\textbf{Statement:}
\begin{equation}
\gamma_{FS} < 0 \quad \text{when } \Psi_S = \text{``prosocial framing''}
\end{equation}

\textbf{Interpretation:} Financial incentives ($u_F$) and social approval ($u_S$)
are often substitutes in prosocial contexts. Paying for blood donations can reduce
the social signaling value of donation.

\textbf{Empirical Basis:} \citet{gneezy2011pay}, \citet{ariely2009doing}
\end{tcolorbox}

% UNMAPPED_CMP-7
\begin{tcolorbox}[colback=gray!5!white, colframe=gray!75!black,
    title=Axiom UNMAPPED_CMP-7: Positive Definiteness Requirement]
\textbf{Statement:}
\begin{equation}
\mathbf{\Gamma} \text{ must be positive semi-definite for well-defined utility}
\end{equation}

\textbf{Interpretation:} The γ-matrix must satisfy eigenvalue conditions to ensure
utility is well-defined (no explosive interactions).

\textbf{Implication:} Not all combinations of γ values are permissible.
Strong mutual crowding-out across many dimensions is constrained.
\end{tcolorbox}

% UNMAPPED_CMP-8
\begin{tcolorbox}[colback=gray!5!white, colframe=gray!75!black,
    title=Axiom UNMAPPED_CMP-8: Emergence from Complementarity]
\textbf{Statement:}
\begin{equation}
\gamma_{ij} > 0 \Rightarrow \frac{\partial K}{\partial Q} \text{ can be non-monotonic}
\end{equation}

\textbf{Interpretation:} When components are complements, the marginal value of
quantity $Q$ can increase with quality $K$. This creates the ``non-concavity''
that enables emergence phenomena.

\textbf{Cross-Reference:} Chapter 7 details how positive γ enables optimal bundling.
\end{tcolorbox}

% UNMAPPED_CMP-9
\begin{tcolorbox}[colback=gray!5!white, colframe=gray!75!black,
    title=Axiom UNMAPPED_CMP-9: Level-Level Complementarity]
\textbf{Statement:}
\begin{equation}
\gamma_{LL'} = \frac{\partial^2 W}{\partial U_L \partial U_{L'}} \quad \text{for welfare levels } L, L'
\end{equation}

\textbf{Interpretation:} Complementarity applies not just between utility dimensions
but between welfare levels. Individual and organizational utility can be complements.

\textbf{Cross-Reference:} Level structure from CORE-WHO (AAA).
\end{tcolorbox}

% UNMAPPED_CMP-10
\begin{tcolorbox}[colback=gray!5!white, colframe=gray!75!black,
    title=Axiom UNMAPPED_CMP-10: Intervention Complementarity]
\textbf{Statement:}
\begin{equation}
\gamma_{I_a I_b} = \frac{\partial^2 E[\Delta U]}{\partial I_a \partial I_b}
\end{equation}

\textbf{Interpretation:} Interventions themselves have complementarity structure.
Information campaigns ($I_a$) may complement or substitute for incentive changes ($I_b$).

\textbf{Cross-Reference:} Full treatment in CORE-EIT (IE).
\end{tcolorbox}

% UNMAPPED_CMP-11
\begin{tcolorbox}[colback=gray!5!white, colframe=gray!75!black,
    title=Axiom UNMAPPED_CMP-11: Narrative Modulation of Complementarity]
\textbf{Statement:}
\begin{equation}
\gamma_{ij}(N) = \gamma_{ij}^{\text{base}} + \Delta\gamma(N) \quad \text{where } \text{sgn}(\gamma_{ij}(N)) \text{ can differ from } \text{sgn}(\gamma_{ij}^{\text{base}})
\end{equation}

\textbf{Interpretation:} Prevailing narratives can \emph{reverse} the sign of complementarity.
This is stronger than UNMAPPED_CMP-4 (context dependence): narratives don't just modulate magnitude,
they can transform complements into substitutes and vice versa.

\textbf{Mechanism:}
\begin{itemize}[nosep]
\item \textbf{Exculpatory narrative} $N_E$: ``Others are responsible'' $\Rightarrow$ $\gamma < 0$ (substitutes: their action replaces mine)
\item \textbf{Responsibilizing narrative} $N_R$: ``We are all responsible'' $\Rightarrow$ $\gamma > 0$ (complements: we must act together)
\end{itemize}

\textbf{Formal Statement:}
\begin{equation}
\gamma_{ij}(N) = \begin{cases}
\gamma^{+} > 0 & \text{if } N = N_R \text{ (responsibilizing)} \\
\gamma^{-} < 0 & \text{if } N = N_E \text{ (exculpatory)}
\end{cases}
\end{equation}

\textbf{Empirical Basis:} \citet{benabou2018narratives} --- moral narratives as strategic complements/substitutes.

\textbf{Cross-Reference:} Full formalization in Appendix UNMAPPED_RB (LIT-BENABOU), Axiom UNMAPPED_RB-7.
\end{tcolorbox}


%%%%%%%%%%%%%%%%%%%%%%%%%%%%%%%%%%%%%%%%%%%%%%%%%%%%%%%%%%%%%%%%%%%%%%%%%%%%%%%
\section{Cross-Dimension Complementarity}
\label{sec:how-dimension}
%%%%%%%%%%%%%%%%%%%%%%%%%%%%%%%%%%%%%%%%%%%%%%%%%%%%%%%%%%%%%%%%%%%%%%%%%%%%%%%

\subsection{The FEPSDE γ-Matrix}

Applying complementarity to the utility dimensions from CORE-WHAT:

\begin{equation}
\mathbf{\Gamma}_{\text{FEPSDE}} = \begin{pmatrix}
1 & \gamma_{FE} & \gamma_{FP} & \gamma_{FS} & \gamma_{FD} & \gamma_{FX} \\
\gamma_{EF} & 1 & \gamma_{EP} & \gamma_{ES} & \gamma_{ED} & \gamma_{EX} \\
\gamma_{PF} & \gamma_{PE} & 1 & \gamma_{PS} & \gamma_{PD} & \gamma_{PX} \\
\gamma_{SF} & \gamma_{SE} & \gamma_{SP} & 1 & \gamma_{SD} & \gamma_{SX} \\
\gamma_{DF} & \gamma_{DE} & \gamma_{DP} & \gamma_{DS} & 1 & \gamma_{DX} \\
\gamma_{XF} & \gamma_{XE} & \gamma_{XP} & \gamma_{XS} & \gamma_{XD} & 1
\end{pmatrix}
\end{equation}

where:
\begin{itemize}[nosep]
\item F = Financial, E = Experiential, P = Psychological
\item S = Social, D = Developmental, X = Existential
\end{itemize}

\subsection{Known Interaction Patterns}

\begin{table}[htbp]
\centering
\caption{Empirically Established γ Patterns}
\label{tab:how-empirical-gamma}
\renewcommand{\arraystretch}{1.3}
\begin{tabular}{@{}cclc@{}}
\toprule
\textbf{Pair} & \textbf{Typical γ} & \textbf{Pattern} & \textbf{Source} \\
\midrule
F $\leftrightarrow$ X & $-0.3$ & Crowding-out & \citet{frey1997relationship} \\
F $\leftrightarrow$ S & $-0.2$ & Social signaling reduction & \citet{gneezy2011pay} \\
P $\leftrightarrow$ S & $+0.4$ & Autonomy-relatedness synergy & \citet{ryan2000self} \\
S $\leftrightarrow$ D & $+0.3$ & Social learning effects & \citet{bandura1977social} \\
E $\leftrightarrow$ X & $+0.5$ & Experience-meaning synergy & \citet{kahneman1999wellbeing} \\
\bottomrule
\end{tabular}
\end{table}


%%%%%%%%%%%%%%%%%%%%%%%%%%%%%%%%%%%%%%%%%%%%%%%%%%%%%%%%%%%%%%%%%%%%%%%%%%%%%%%
\section{Cross-Level Complementarity}
\label{sec:how-level}
%%%%%%%%%%%%%%%%%%%%%%%%%%%%%%%%%%%%%%%%%%%%%%%%%%%%%%%%%%%%%%%%%%%%%%%%%%%%%%%

\subsection{The Level γ-Matrix}

Applying complementarity across welfare levels from CORE-WHO:

\begin{equation}
\mathbf{\Gamma}_{\text{Level}} = \begin{pmatrix}
1 & \gamma_{01} & \gamma_{02} & \cdots \\
\gamma_{10} & 1 & \gamma_{12} & \cdots \\
\gamma_{20} & \gamma_{21} & 1 & \cdots \\
\vdots & \vdots & \vdots & \ddots
\end{pmatrix}
\end{equation}

\subsection{Typical Level Interactions}

\begin{itemize}
\item \textbf{Individual-Organizational} ($\gamma_{03}$): Often positive --- when
      organizations thrive, individuals within them benefit (enabling effect)
\item \textbf{Individual-Societal} ($\gamma_{0n}$): Can be negative --- individual
      optimization may create negative externalities for society
\item \textbf{Team-Organizational} ($\gamma_{23}$): Usually positive --- strong
      teams create strong organizations
\end{itemize}


%%%%%%%%%%%%%%%%%%%%%%%%%%%%%%%%%%%%%%%%%%%%%%%%%%%%%%%%%%%%%%%%%%%%%%%%%%%%%%%
\section{Integration with Other CORE Appendices}
\label{sec:how-integration}
%%%%%%%%%%%%%%%%%%%%%%%%%%%%%%%%%%%%%%%%%%%%%%%%%%%%%%%%%%%%%%%%%%%%%%%%%%%%%%%

\begin{tcolorbox}[colback=yellow!5!white, colframe=yellow!75!black,
    title=Integration: CORE-HOW $\leftrightarrow$ CORE-WHO]
\textbf{What HOW provides to WHO:}
\begin{itemize}[nosep]
\item Cross-level complementarity values $\gamma_{LL'}$
\item Determines whether levels enhance or crowd each other
\end{itemize}

\textbf{What HOW requires from WHO:}
\begin{itemize}[nosep]
\item Level definitions ($L_0, \ldots, L_n$)
\item Entity structure at each level
\end{itemize}
\end{tcolorbox}

\begin{tcolorbox}[colback=yellow!5!white, colframe=yellow!75!black,
    title=Integration: CORE-HOW $\leftrightarrow$ CORE-WHAT]
\textbf{What HOW provides to WHAT:}
\begin{itemize}[nosep]
\item Cross-dimension complementarity values $\gamma_{ij}$
\item Determines which dimension combinations are synergistic
\end{itemize}

\textbf{What HOW requires from WHAT:}
\begin{itemize}[nosep]
\item Dimension definitions (FEPSDE taxonomy)
\item Measurement scales for each dimension
\end{itemize}
\end{tcolorbox}

\begin{tcolorbox}[colback=yellow!5!white, colframe=yellow!75!black,
    title=Integration: CORE-HOW $\leftrightarrow$ CORE-WHEN]
\textbf{What HOW provides to WHEN:}
\begin{itemize}[nosep]
\item Baseline γ values to be modulated by context
\item Interaction structure that context affects
\end{itemize}

\textbf{What HOW requires from WHEN:}
\begin{itemize}[nosep]
\item Context specification $\Psi$
\item Rules for how $\gamma_{ij}(\Psi)$ varies with context
\end{itemize}
\end{tcolorbox}

\begin{tcolorbox}[colback=green!5!white, colframe=green!75!black,
    title=Application Domain: Complex Problem Solving (CP)]
\textbf{CPS leverages γ as ``Vernetztheit'' (Interconnectedness):}
\begin{itemize}[nosep]
\item Dörner's second quality: Problem variables influence each other
\item In CPS: $\gamma_{ij} > 0$ means changing $X_i$ affects $X_j$
\item High connectivity increases problem complexity exponentially
\end{itemize}

\textbf{Key Finding:} Variable interconnectedness ($\gamma$) is the second-strongest lever ($\sim 28\%$ influence) in CPS after intransparency.

\textbf{Cross-Reference:} See Appendix UNMAPPED_CP (DOMAIN-CPS) for the complete 10C-CPS mapping.
\end{tcolorbox}


%%%%%%%%%%%%%%%%%%%%%%%%%%%%%%%%%%%%%%%%%%%%%%%%%%%%%%%%%%%%%%%%%%%%%%%%%%%%%%%
\section{Worked Example}
\label{sec:how-example}
%%%%%%%%%%%%%%%%%%%%%%%%%%%%%%%%%%%%%%%%%%%%%%%%%%%%%%%%%%%%%%%%%%%%%%%%%%%%%%%

\subsection{Example: Intervention Portfolio Design}

\paragraph{Setup.}
A company wants to increase employee motivation through:
\begin{itemize}
\item Financial bonus ($I_F$)
\item Social recognition program ($I_S$)
\item Autonomy enhancement ($I_P$)
\end{itemize}

\paragraph{Application.}
Using the γ-matrix from Table~\ref{tab:how-empirical-gamma}:

\begin{align}
\Delta U &= I_F \cdot \Delta u_F + I_S \cdot \Delta u_S + I_P \cdot \Delta u_P \\
&\quad + \gamma_{FS} \cdot I_F \cdot I_S + \gamma_{FP} \cdot I_F \cdot I_P + \gamma_{SP} \cdot I_S \cdot I_P \\
&= I_F \cdot \Delta u_F + I_S \cdot \Delta u_S + I_P \cdot \Delta u_P \\
&\quad - 0.2 \cdot I_F \cdot I_S + 0.1 \cdot I_F \cdot I_P + 0.4 \cdot I_S \cdot I_P
\end{align}

\paragraph{Result.}
The negative $\gamma_{FS} = -0.2$ means combining financial bonus with social recognition
reduces effectiveness (financial framing undermines social signaling).
The positive $\gamma_{SP} = +0.4$ means combining social recognition with autonomy is synergistic.

\paragraph{Recommendation.}
Don't combine $I_F$ and $I_S$ simultaneously. Either:
\begin{enumerate}
\item Use $(I_S + I_P)$ without $I_F$ for intrinsically motivated employees
\item Use $I_F$ alone for extrinsically motivated employees
\end{enumerate}


%%%%%%%%%%%%%%%%%%%%%%%%%%%%%%%%%%%%%%%%%%%%%%%%%%%%%%%%%%%%%%%%%%%%%%%%%%%%%%%
\section{Implications for Practice}
\label{sec:how-practice}
%%%%%%%%%%%%%%%%%%%%%%%%%%%%%%%%%%%%%%%%%%%%%%%%%%%%%%%%%%%%%%%%%%%%%%%%%%%%%%%

\begin{tcolorbox}[colback=green!5!white, colframe=green!75!black,
    title=Practical Implications]
\begin{enumerate}
\item \textbf{Never Assume Independence:} Before combining interventions or
      benefits, check the γ-matrix for crowding-out risks
\item \textbf{Exploit Synergies:} When $\gamma > 0$, bundle components together
      for effects larger than the sum
\item \textbf{Sequence Strategically:} If $\gamma < 0$, consider sequential
      rather than simultaneous implementation
\item \textbf{Context Matters:} The same components may be complements or
      substitutes depending on framing and context
\end{enumerate}
\end{tcolorbox}


%%%%%%%%%%%%%%%%%%%%%%%%%%%%%%%%%%%%%%%%%%%%%%%%%%%%%%%%%%%%%%%%%%%%%%%%%%%%%%%
\section{Summary}
\label{sec:how-summary}
%%%%%%%%%%%%%%%%%%%%%%%%%%%%%%%%%%%%%%%%%%%%%%%%%%%%%%%%%%%%%%%%%%%%%%%%%%%%%%%

\begin{tcolorbox}[colback=gray!5!white, colframe=gray!75!black,
    title=Appendix UNMAPPED_B Summary: CORE-HOW]

\textbf{Core Contribution:}
Establishes that utility components interact through the complementarity parameter
$\gamma_{ij} \in [-1, +1]$, enabling analysis of synergies and crowding-out.

\textbf{Key Formula:}
\begin{equation}
U = \sum_i u_i + \sum_{i < j} \gamma_{ij} \cdot u_i \cdot u_j
\end{equation}

\textbf{Key Insights:}
\begin{itemize}[nosep]
\item Additive separability ($\gamma = 0$) is a special case, not the general case
\item Financial and intrinsic motivation typically crowd out ($\gamma_{FX} < 0$)
\item Complementarity enables emergence and non-concave utility
\item Intervention portfolios must account for cross-intervention γ
\end{itemize}

\textbf{Next Steps:}
See Appendix UNMAPPED_C (CORE-WHAT) for the dimensions that interact.
See Appendix UNMAPPED_V (CORE-WHEN) for context-dependent γ values.

\end{tcolorbox}


%%%%%%%%%%%%%%%%%%%%%%%%%%%%%%%%%%%%%%%%%%%%%%%%%%%%%%%%%%%%%%%%%%%%%%%%%%%%%%%
\section{Glossary of Symbols}
\label{sec:how-glossary}
%%%%%%%%%%%%%%%%%%%%%%%%%%%%%%%%%%%%%%%%%%%%%%%%%%%%%%%%%%%%%%%%%%%%%%%%%%%%%%%

\begin{table}[htbp]
\centering
\caption{Symbols Introduced in Appendix UNMAPPED_B}
\label{tab:how-symbols}
\renewcommand{\arraystretch}{1.3}
\begin{tabular}{@{}clp{6cm}c@{}}
\toprule
\textbf{Symbol} & \textbf{Name} & \textbf{Definition} & \textbf{Status} \\
\midrule
$\gamma_{ij}$ & Complementarity & Cross-partial derivative of utility & NEW \\
$\mathbf{\Gamma}$ & γ-Matrix & Complete interaction specification & NEW \\
$\gamma_{ij}(N)$ & Narrative-modulated γ & Complementarity conditional on narrative $N$ & NEW \\
$N_E, N_R$ & Narrative types & Exculpatory vs.\ responsibilizing narratives & NEW \\
FEPSDE & Dimensions & Financial, Experiential, Psychological, Social, Developmental, Existential & \checkmark \\
\bottomrule
\end{tabular}
\end{table}


%%%%%%%%%%%%%%%%%%%%%%%%%%%%%%%%%%%%%%%%%%%%%%%%%%%%%%%%%%%%%%%%%%%%%%%%%%%%%%%
\section{Critical Foundations and Objections}
\label{sec:how-critical}
%%%%%%%%%%%%%%%%%%%%%%%%%%%%%%%%%%%%%%%%%%%%%%%%%%%%%%%%%%%%%%%%%%%%%%%%%%%%%%%

\subsection{Objection 1: Measurement Problem}

\begin{tcolorbox}[colback=red!5!white, colframe=red!50!black,
    title=Objection: γ Cannot Be Measured]
\textit{``How can we measure interaction effects when we can barely measure utility itself?
The γ-matrix is theoretically elegant but empirically empty.''}
\end{tcolorbox}

\textbf{Response:}
γ is identifiable from factorial experimental designs that vary two components
simultaneously. The interaction term in $2 \times 2$ designs directly estimates γ.
Meta-analyses (e.g., \citet{deci1999meta}) provide robust γ estimates for key pairs.

\subsection{Objection 2: Complexity Explosion}

\begin{tcolorbox}[colback=red!5!white, colframe=red!50!black,
    title=Objection: Too Many Parameters]
\textit{``With $n$ components, the γ-matrix has $n(n-1)/2$ parameters.
This is overparameterized and unfalsifiable.''}
\end{tcolorbox}

\textbf{Response:}
Not all γ values need separate estimation. Theoretical constraints (UNMAPPED_CMP-5, UNMAPPED_CMP-6)
provide many values a priori. Empirically, most γ values cluster around zero;
we need only estimate the non-zero interactions.


%%%%%%%%%%%%%%%%%%%%%%%%%%%%%%%%%%%%%%%%%%%%%%%%%%%%%%%%%%%%%%%%%%%%%%%%%%%%%%%
\section{Open Issues and Research Agenda}
\label{sec:how-open}
%%%%%%%%%%%%%%%%%%%%%%%%%%%%%%%%%%%%%%%%%%%%%%%%%%%%%%%%%%%%%%%%%%%%%%%%%%%%%%%

\begin{table}[htbp]
\centering
\caption{Research Roadmap for CORE-HOW}
\label{tab:how-roadmap}
\renewcommand{\arraystretch}{1.3}
\begin{tabular}{@{}clccl@{}}
\toprule
\textbf{\#} & \textbf{Issue} & \textbf{Priority} & \textbf{Status} & \textbf{Requirements} \\
\midrule
1 & Complete FEPSDE γ-matrix calibration & HIGH & Partial & Meta-analysis of factorial studies \\
2 & Context modulation rules for γ & HIGH & Open & Moderator analysis \\
3 & Higher-order interactions (γ₃) & MEDIUM & Open & Three-way factorial designs \\
4 & Dynamic γ (time-varying) & MEDIUM & Open & Panel data methods \\
5 & Neural basis of complementarity & LOW & Open & Neuroscience collaboration \\
\bottomrule
\end{tabular}
\end{table}


%%%%%%%%%%%%%%%%%%%%%%%%%%%%%%%%%%%%%%%%%%%%%%%%%%%%%%%%%%%%%%%%%%%%%%%%%%%%%%%
\section{References}
\label{sec:how-references}
%%%%%%%%%%%%%%%%%%%%%%%%%%%%%%%%%%%%%%%%%%%%%%%%%%%%%%%%%%%%%%%%%%%%%%%%%%%%%%%

\begin{tcolorbox}[colback=blue!5!white, colframe=blue!75!black]
\textbf{Master Bibliography:} See \texttt{bcm\_master.bib}
\end{tcolorbox}

\subsection{Key References}

\begin{itemize}[nosep]
\item \textbf{Crowding-Out:} \citet{frey1997relationship}, \citet{deci1999meta}
\item \textbf{Social Incentives:} \citet{gneezy2011pay}, \citet{ariely2009doing}
\item \textbf{Self-Determination:} \citet{ryan2000self}
\item \textbf{Complementarity Theory:} \citet{milgrom1990job}
\end{itemize}

\nocite{bcm_master}

%%%%%%%%%%%%%%%%%%%%%%%%%%%%%%%%%%%%%%%%%%%%%%%%%%%%%%%%%%%%%%%%%%%%%%%%%%%%%%%
% FOOTER
%%%%%%%%%%%%%%%%%%%%%%%%%%%%%%%%%%%%%%%%%%%%%%%%%%%%%%%%%%%%%%%%%%%%%%%%%%%%%%%

\vspace{1em}
\hrule
\vspace{0.5em}

\noindent\textit{Cross-references:}
Appendix UNMAPPED_GLS (Glossary),
Appendix UNMAPPED_AAA (CORE-WHO),
Appendix UNMAPPED_C (CORE-WHAT),
Appendix UNMAPPED_V (CORE-WHEN),
Appendix UNMAPPED_IE (CORE-EIT),
Appendix UNMAPPED_CAL (METHOD-LLMMC),
Appendix UNMAPPED_RB (LIT-BENABOU) --- Axiom UNMAPPED_CMP-11 narrative modulation.

\noindent\textit{Master Bibliography:} \texttt{bcm\_master.bib}

% =============================================================================
% END OF APPENDIX B
% =============================================================================
